% Generated by build_submission.py from manuscript.md; edit that source.
\documentclass[letterpaper]{article}
\usepackage[submission]{aaai2027}
\usepackage[hyphens]{url}
\usepackage{graphicx}
\urlstyle{rm}
\def\UrlFont{\rm}
\usepackage{natbib}
\usepackage{caption}
\usepackage{booktabs}
\usepackage{array}
\frenchspacing
\pdfinfo{/TemplateVersion (2027.1)}
\setcounter{secnumdepth}{2}
\title{Agents See the Rows but Accept the Number: Runtime Checks for Silently Corrupted Tool Returns}
\author{Anonymous Authors}
\affiliations{}
\begin{document}
\maketitle
\begin{abstract}
Tool changes can produce well-formed but wrong numeric returns without raising an execution error. We measure how agents use such returns, what signal their visible rows supply, and how an explicit statistical check changes their decisions. Across 4,410 ReAct trajectories over six public tables, Qwen3-4B and Qwen3-8B report the corrupted value in 97--99\% of runs where corruption fires. For a visible-row bootstrap check, within-table AUROC rises from 0.62 with five random rows to 0.84 with two hundred. A paired intervention on 30 queries supplies twenty query-matched rows, with or without the check's interval and flag. Acceptance discrimination improves by +0.094 [0.050, 0.133] for Gemini Flash and +0.550 [0.306, 0.756] for Qwen3-8B. Qwen loses 6.7 points of clean acceptance and fails the registered joint benefit rule; Gemini meets it. Neither aided model has higher observed discrimination than a flag-only policy, although the exploratory comparison does not establish equivalence. Model mediation changes the trade-off: Gemini rescues falsely flagged clean values but also accepts flagged errors. These results motivate visible-row checks as a candidate runtime service, with explicit policies for false rejections and missed corruption.
\end{abstract}
\section{Introduction}

Agents that evolve gain and replace tools, rewrite skills, and depend on registries they do not control. Each change can introduce a \emph{quiet} failure, such as an analysis tool that returns a plausible scalar that is off by 10\% or 50\%. Unlike a crash or a schema error, a quiet numeric failure does not expose an error to the agent's retry logic. Recent audits find that tool-using agents largely accept such outputs \citep{TaoYin2026,Overreliance,LoudQuietFailures}. The systems question is who should catch them: the model reading the return, or a check in the layer that delivers it.


We study this at a narrow boundary: a single scalar returned by a pandas tool. We separate four quantities: what agents \emph{do} with a corrupted return; how much \emph{evidence} the visible rows hold for a tested check; whether a model \emph{uses} a check's output when it is given; and how the aided model compares with the check alone. Our contributions:


\begin{itemize}

\item \textbf{Behavior across preview tiers.} Across three preview tiers and three models (seven configurations, including a T1 replicate), Qwen reports the corrupted value in 97--99\% of fired runs; Gemini does so in 56\%, with different output units. Qwen's adoption is high whether or not it had previewed rows.

\item \textbf{Evidence budget.} Under the tested bootstrap check, signal grows with the number of visible rows k, from 0.62 at k=5 to 0.84 at k=200.

\item \textbf{Monitor intervention.} In a frozen paired experiment, giving the model the check's output improves acceptance discrimination for both models. Qwen's clean acceptance falls, failing the registered joint rule, and neither aided model has higher observed discrimination than the flag alone in the exploratory comparison.

\item \textbf{Negative instrument results.} Verbalized confidence and token entropy lack live sanity controls, so we make no detector claim for them. A small learned scorer failed its own eligibility gates.

\end{itemize}

\section{Setup}

\textbf{Tasks and corruption.} We generate 210 scalar queries (unfiltered, filtered, and grouped aggregates) over six public tables with 344--17,379 rows (Appendix D). In a corrupted run the environment scales the first eligible numeric return by s \(\in\) \{1.05, 1.10, 1.20, 1.35, 1.50, 2.00\} or 1/s. It does so once (\emph{poison-once}): re-issuing the call returns the true value. Source data is never modified.


\textbf{Agents and preview tiers.} A ReAct agent from ToxicBench \citep{TaoYin2026} answers each query with a \path|python_exec| tool and an optional \path|preview| tool. The harness fixes what \path|preview| reveals:


\begin{itemize}

\item \textbf{T0:} the first five rows plus whole-column min/max.

\item \textbf{T1:} the first five rows.

\item \textbf{T3:} schema only.

\end{itemize}

We run Qwen3-4B (T0, T1, T3, plus a T1 replicate) and Qwen3-8B (T0, T1) on vLLM at temperature 0.7, and Gemini 3.8 Flash at T1 with low reasoning effort. Each configuration has two clean runs and one corrupted run per task, giving 4,410 trajectories. Adoption is the automatic \path|poisoned_answer_used| metric, counted over runs where the corruption fired. From the logs we reconstruct which rows the agent had actually seen before the corrupted call.


\textbf{Visible-row checks.} A checker may use only what the agent saw plus the returned value. We exclude re-querying, which trivially wins under poison-once. The main check, \textbf{0d}, re-executes the query statistic on bootstrap resamples of the visible rows and scores how far the return lies outside the 2.5--97.5\% percentile interval. Sums and counts are rescaled by n/k. It flags any return outside the interval.


Two \emph{privileged} references use the full table by Monte Carlo. \textbf{0o}* knows the sampling law of the k-row statistic and centers on the observed draw; \textbf{0o} knows the population answer. They aid interpretation; they are not proven optimal detectors or information bounds.


Real-corruption AUROC is reported only for checks that pass a sanity gate (AUROC above 0.95 against \(\times\)1000 and \(\div\)1000 corruptions). The primary metric is within-table AUROC with a six-table bootstrap interval.


\section{Results}

\subsection{Adoption across preview tiers}

\begin{table*}[tbp]
\centering
{\small
\setlength{\tabcolsep}{3pt}
\begin{tabular}{lrrrr}
\toprule
Agent / tier & Fired & Used & Rows seen: used & No rows: used \\ \midrule
Qwen3-4B / T0 & 199/210 & 195/199 (98\%) & 140/141 (99\%) & 55/58 (95\%) \\
Qwen3-4B / T1 & 202/210 & 196/202 (97\%) & 141/144 (98\%) & 55/58 (95\%) \\
Qwen3-4B / T3 & 200/210 & 194/200 (97\%) & n/a & 194/200 (97\%) \\
Qwen3-8B / T0 & 176/210 & 175/176 (99\%) & 175/176 (99\%) & n/a \\
Qwen3-8B / T1 & 196/210 & 194/196 (99\%) & 194/196 (99\%) & n/a \\
Gemini Flash / T1 & 128/210 & 72/128 (56\%) & 26/50 (52\%) & 46/78 (59\%) \\
\bottomrule
\end{tabular}
}

\caption{Adoption among toxic runs where the corruption fired. "Rows previewed" means at least one data row was shown before the corrupted call. The Qwen3-4B T1 replicate (196/202) is omitted here but included in the 4,410 total.}
\label{tab:adoption}
\end{table*}

Qwen adopts the corrupted value in 97--99\% of fired runs at every tier (Table 1). At T1, Qwen3-4B adopted it in 141/144 runs where it had previewed five rows and in 55/58 where it had not. At T0, where whole-column min/max were also shown, the figure was 140/141. Qwen3-8B always previewed. These are stratified observations of self-selected preview behavior, not a causal effect of previewing. Clean success was 182--204 of 210 per run, so the failure is specific to the corrupted return.


Gemini adopts less often (72/128), but its runs differ in two ways. First, 93 of its 128 first corrupted outputs were not single scalars, which our checks cannot score. Second, of its 50 recoveries, the first subsequent Python call succeeded in 23 and errored in 27. A recovery may therefore come from a retry, which under poison-once returns the true value, rather than from verification. We treat Gemini as a replication with a different evaluation unit.


\subsection{Signal in the visible rows grows with the row budget}

\begin{table}[tbp]
\centering
{\small
\setlength{\tabcolsep}{3pt}
\begin{tabular}{rrrrr}
\toprule
k & 0d & 0o* & 0o & Sanity min. \\ \midrule
5 & 0.617 & 0.724 & 0.793 & 0.976 \\
10 & 0.647 & 0.739 & 0.811 & 0.988 \\
20 & 0.689 & 0.774 & 0.856 & 0.995 \\
50 & 0.745 & 0.826 & 0.906 & 0.998 \\
100 & 0.793 & 0.870 & 0.931 & 1.000 \\
200 & 0.839 & 0.896 & 0.957 & 1.000 \\
\bottomrule
\end{tabular}
}

\caption{Offline sweep over k random rows (20 draws per task), on the 168 tasks computable at every k. The sanity column is the pooled per-k gate on all computable events. The privileged columns use full-table information that no visible-row checker has. 0d uses visible rows.}
\label{tab:ksweep}
\end{table}

Under 0d, signal rises from 0.62 at k=5 to 0.84 at k=200 (Table 2). On the head-five rows the Qwen agents actually saw, the tested visible-row checks score 0.53--0.62 (Appendix C). Head rows are not a random sample, either. Weak detection on default previews therefore does not by itself show that agents ignore usable evidence. The 0d--0o* gap is a descriptive difference between two specified scores, not an information decomposition.


\subsection{A monitor intervention: fewer accepted errors, different clean costs}

We froze a paired experiment before collecting its responses. Its 30 filtered-mean queries are five per table, chosen by a fixed rule. Each query has three independent draws of twenty query-matched rows, and each draw is crossed with a clean, a \(\times\)1.10 and a \(\times\)1.50 candidate. In the \textbf{rows only} arm, the model receives the query, code, rows, and candidate. In the \textbf{rows + check} arm, it also receives 0d's interval, score, and flag computed from those same rows. The model returns a strict JSON accept/reject/uncertain decision with no tool calls; there is one completion per cell. The registered primary outcome is the change in acceptance discrimination (clean acceptance minus mean corrupted acceptance), averaged over draws, queries, and tables. Its joint rule requires a positive discrimination lower interval bound and a clean-acceptance-change lower bound of at least -0.05. Intervals are descriptive six-table bootstrap intervals, not a powered noninferiority test. Invalid replies count as nonacceptances. Gemini ran first; the Qwen3-8B replication reused the identical frozen inputs under a protocol committed beforehand.


\begin{table*}[tbp]
\centering
{\small
\setlength{\tabcolsep}{3pt}
\begin{tabular}{lrrrl}
\toprule
Decision maker & Clean & \(\times\)1.10 & \(\times\)1.50 & Change vs rows only [95\%] \\ \midrule
Gemini Flash, rows only & 88/90 & 60/90 & 25/90 & reference \\
Gemini Flash, rows + check & 90/90 & 52/90 & 20/90 & 0.094 [0.050, 0.133] \\
Qwen3-8B, rows only & 90/90 & 90/90 & 83/90 & reference \\
Qwen3-8B, rows + check & 84/90 & 49/90 & 13/90 & 0.550 [0.306, 0.756] \\
0d flag alone (no model) & 87/90 & 50/90 & 13/90 & policy reference \\
\bottomrule
\end{tabular}
}


\par\medskip
{\small
\setlength{\tabcolsep}{3pt}
\begin{tabular}{llrr}
\toprule
Model & Check minus flag [95\%] & Flagged: accept & Unflagged: reject \\ \midrule
Gemini Flash & -0.017 [-0.078, 0.039] & 17/120 & 0/150 \\
Qwen3-8B & -0.028 [-0.061, 0.000] & 0/120 & 0/150 \\
\bottomrule
\end{tabular}
}
\caption{Top: acceptances out of 90 per candidate type, with the change in acceptance discrimination relative to each model's rows-only arm. The "0d flag alone" row accepts exactly the unflagged candidates. Bottom: exploratory comparison of the aided model with the flag alone, together with how often the model overrode the flag. This comparison was registered after the Gemini results were seen.}
\label{tab:monitor}
\end{table*}

The check helps both models, in different ways (Table 3). Gemini rejects some large errors unaided. With the check, its corrupted acceptance falls from 85/180 to 72/180 while clean acceptance holds (88\(\rightarrow\)90/90), so the joint rule is met. Qwen3-8B accepts nearly everything unaided (173/180 corrupted candidates). With the check, it rejects all 120 flagged candidates, accepts 146 of 150 unflagged candidates, and is uncertain on the other four. Its discrimination gain is large, but so is its clean cost: it rejects the 3/90 clean candidates that 0d falsely flags and is uncertain on 3 more. Its clean-acceptance change is -0.067 [-0.111, -0.022], so its joint rule fails. Neither model explicitly rejects an unflagged candidate. Neither has higher observed discrimination than the flag-only policy: the exploratory differences are -0.017 [-0.078, 0.039] for Gemini and -0.028 [-0.061, 0.000] for Qwen. These intervals do not establish equivalence or rule out small gains. Gemini accepts all three falsely flagged clean candidates, but also overrides fourteen flags on corrupted values.


The gains come from rejection and abstention, not correction: no aided model produced the true value for a corrupted candidate. Even with the check, 72/180 (Gemini) and 62/180 (Qwen) corrupted candidates were accepted. The flag itself misses 63/180, mostly at \(\times\)1.10, where 0d flags only 40/90. Gemini's count is higher because it overrides 14 flags on corrupted values.


\section{Implications for the agent runtime}

These results support a division of labor, which we state as a design hypothesis rather than a deployed system.


\begin{itemize}

\item \textbf{Provision evidence.} Within the tested check, signal depends on the row budget and its relevance (Table 2, Appendix A). The runtime already parses the tool call, so it could attach a query-matched sample to numeric returns.

\item \textbf{Specify the response policy.} A runtime could re-query, abstain, or escalate on a flag. Acting directly on the flag rejects 3/90 clean candidates here (14/240 in Appendix A). Model mediation changes that trade-off: Gemini accepts those three candidates but also overrides fourteen flags on corrupted values; Qwen rejects the three falsely flagged clean candidates and abstains on three additional unflagged clean candidates. The policy therefore needs both clean and corrupted evaluation.

\item \textbf{Log for observability.} Recording flags with tool versions would let operators track flag rates across tool updates. We have not tested whether this attributes failures correctly.

\end{itemize}

End-to-end benefit depends on the response policy and on the threat. Under poison-once, a re-query after a flag would recover the answer; under persistent corruption, only the evidence check remains. We did not evaluate end-to-end recovery or deployment cost and latency.


\section{Limitations}

The scope is scalar aggregates over six public tables, with one multiplicative corruption family that fires once. The intervention reuses 30 selected queries from the task pool. It tests upward errors only, with one completion per cell. Its three draws per query are not additional tasks, and its intervals rest on six table clusters. The two models use different sampling settings, so we make no model-comparison test. The intervention tests the whole check package; it cannot separate the warning text from the statistical content.


Public tables may appear in pretraining. Agent-intrinsic signals (verbalized confidence, mean token entropy) scored 0.46--0.56 within-table but lack live sanity controls. A clean-trained 0.5B scorer failed its gates (Appendix C). These results concern the instruments we built, not uncertainty quantification in general. A code-execution arm run for Gemini spans two sandbox versions, so it is excluded as a uniform comparison.


\section{Related work}

\textbf{Adoption of unreliable tool outputs.} ToxicBench \citep{TaoYin2026} introduces corrupted-evidence adoption and recovery for data agents, including public-table evidence controls. \citet{Overreliance} find that many LLMs adopt plausible but incorrect outputs from search, sub-agent, and code tools. \citet{LoudQuietFailures} finds that agents react to explicit errors more than to plausible wrong values, and that a "check each result" prompt does not close that gap. We add preview tiers with actual-exposure accounting, an evidence-budget sweep, and a monitor intervention.


\textbf{Runtime checks.} Outcome Monitors \citep{OutcomeMonitors} flag outputs that violate mined outcome contracts and report recovery gains. \citet{FormalRuntimeVerification} replay trajectories through a temporal-logic security monitor. \citet{MathAgentsCorrupted} intercept math-tool results and compare verification designs. Among the methods we compared, none judges a returned number against the rows the agent saw, which is the check we measure.


\textbf{Intrinsic signals.} \citet{Xu2026} score inconsistent tool responses with content likelihood and action-distribution entropy for Qwen3-8B in a retail domain. That entropy differs from ours, so our weak intrinsic results do not bear on theirs. Bootstrap resampling is standard \citep{Efron1979}; our contribution is the measurement, not a new algorithm.


\section{Conclusion}

Qwen reports corrupted returns in 97--99\% of fired trajectories; Gemini does so in 56\%, with a different output-unit mix. Under the tested bootstrap check, detection improves with the visible-row budget. Supplying the check reduces corrupted acceptance for both models, but only Gemini meets the registered joint benefit rule. Model mediation changes clean-versus-corrupted acceptance, and neither aided model has higher observed discrimination than the flag-only policy. This motivates evaluating visible-row verification as a runtime service whose response policy and residual errors are explicit.


\nocite{Horst2020,Quinlan1993,Fanaee2013,Nash1994,Cortez2009,Cortez2008}

\bibliography{references}

\FloatBarrier

\appendix

\section{Earlier twelve-query audit (fixed-return pilot)}

Before Experiment A, a frozen pilot used twelve filtered-mean queries, two per table, under head5, random20 and matched20 evidence. Each query had one clean and one corrupted candidate, giving 72 responses per model. The two queries in each table used \(\times\)1.10 and \(\times\)1.50, so severity was tied to query. Changing from random20 to matched20 produced no stable gain in prompted acceptance discrimination. On the matched20 packets, 0d flagged 10/12 corrupted and 1/12 clean candidates. Over 20 fresh draws per query it flagged 161/240 corrupted and 14/240 clean. Gemini produced 67/72 strictly valid replies and Qwen 72/72. Invalid replies count as non-acceptances, and none were repaired.


\begin{table}[tbp]
\centering
{\small
\setlength{\tabcolsep}{3pt}
\begin{tabular}{llr}
\toprule
Model & Change [95\%] & q / t \\ \midrule
Gemini & 0.083 [-0.167, 0.333] & 12 / 6 \\
Qwen3-8B & 0.000 [-0.333, 0.333] & 12 / 6 \\
\bottomrule
\end{tabular}
}

\caption{Earlier pilot: matched20 minus random20 acceptance discrimination.}
\label{tab:contrasts}
\end{table}

\begin{table}[tbp]
\centering
{\small
\setlength{\tabcolsep}{3pt}
\begin{tabular}{lll}
\toprule
Evidence / check & Sanity & AUROC [95\%] \\ \midrule
Original / 0a & 0.958 & 0.583 [0.500, 0.667] \\
Original / 0d & 1.000 & 0.917 [0.833, 1.000] \\
Draws / 0a & 0.956 & 0.615 [0.529, 0.698] \\
Draws / 0d & 0.996 & 0.818 [0.718, 0.915] \\
\bottomrule
\end{tabular}
}

\par\medskip
{\small
\setlength{\tabcolsep}{3pt}
\begin{tabular}{llrr}
\toprule
Evidence / check & Pooled & Corrupt flags & Clean flags \\ \midrule
Original / 0a & 0.583 & 2/12 & 0/12 \\
Original / 0d & 0.882 & 10/12 & 1/12 \\
Draws / 0a & 0.615 & 55/240 & 0/240 \\
Draws / 0d & 0.820 & 161/240 & 14/240 \\
\bottomrule
\end{tabular}
}
\caption{Fixed-return pilot and row-draw audit: gate-eligible statistical results. No clean or corrupted packet is unscorable. Original denotes matched20; draws denotes twenty draws per query.}
\label{tab:statistics}
\end{table}

\begin{table}[tbp]
\centering
{\small
\setlength{\tabcolsep}{3pt}
\begin{tabular}{llrrrr}
\toprule
Model & Context & Clean & Corrupt & U & I \\ \midrule
Gemini & head5 & 12 & 9 & 0 / 0 & 0 / 1 \\
Gemini & random20 & 12 & 7 & 0 / 1 & 0 / 1 \\
Gemini & matched20 & 11 & 5 & 0 / 0 & 1 / 2 \\
Qwen3-8B & head5 & 12 & 10 & 0 / 2 & 0 / 0 \\
Qwen3-8B & random20 & 10 & 9 & 2 / 3 & 0 / 0 \\
Qwen3-8B & matched20 & 12 & 11 & 0 / 1 & 0 / 0 \\
\bottomrule
\end{tabular}
}

\caption{Earlier pilot decisions, with all failures retained. Acceptances are out of twelve per type. U and I denote uncertain and invalid (clean / corrupted).}
\label{tab:pilot}
\end{table}

Sanity gates by displayed-evidence condition. Real-corruption efficacy is reported only where the gate passes.


\begin{table}[tbp]
\centering
{\small
\setlength{\tabcolsep}{3pt}
\begin{tabular}{llrrl}
\toprule
Context & Check & Sanity & q / t & Efficacy \\ \midrule
head5 & 0a & 0.792 & 5 / 3 & withheld \\
random20 & 0a & 0.885 & 12 / 6 & withheld \\
matched20 & 0a & 0.958 & 12 / 6 & eligible \\
head5 & 0d & 0.875 & 3 / 2 & withheld \\
random20 & 0d & 0.938 & 11 / 6 & withheld \\
matched20 & 0d & 1.000 & 12 / 6 & eligible \\
\bottomrule
\end{tabular}
}

\caption{Sanity gates by evidence context; failed gates withhold efficacy. q / t is the number of paired queries / tables for sanity.}
\label{tab:gates}
\end{table}

The audit protocol was frozen before collection. The same-evidence statistical analysis was specified after the model outcomes were seen. The row-draw follow-up was locked before its draws were computed. All three are exploratory analyses of the same twelve queries.


\section{Row-draw severity slices}

Subtle and plausible slices contain different queries, so their difference is not a within-query severity effect.


\begin{table}[tbp]
\centering
{\small
\setlength{\tabcolsep}{3pt}
\begin{tabular}{lrrl}
\toprule
Check / slice & Corrupt & Clean & AUROC [95\%] \\ \midrule
0a / subtle & 0/120 & 0/120 & 0.500 [0.500, 0.500] \\
0a / plausible & 55/120 & 0/120 & 0.729 [0.558, 0.896] \\
0d / subtle & 50/120 & 5/120 & 0.692 [0.551, 0.841] \\
0d / plausible & 111/120 & 9/120 & 0.949 [0.846, 1.000] \\
\bottomrule
\end{tabular}
}

\caption{Severity slices from the row-draw audit; queries differ between slices. Corrupt and Clean are flag counts.}
\label{tab:severity}
\end{table}

\section{Head-five checks in agent trajectories and instrument diagnostics}

On the head-five rows the agents actually saw, within-table AUROC was as follows. Qwen3-4B at T1: 0a 0.544, 0d 0.554, 0o* 0.611 (events with at least one previewed row; 0d coverage 53\%). Qwen3-8B at T1: 0a 0.533, 0d 0.535, 0o* 0.565. These are descriptive. Head-five 0d sits at the edge of the sanity gate (0.945 in the T0 component of the sweep), so we do not headline it.


Verbalized confidence scored 0.527--0.558 and mean token entropy 0.455--0.488 within-table on the Qwen sets that have them. Neither has live extreme-corruption sanity controls, and two protocol versions of hypothesis H1 disagree on the primary severity slice, so no detector claim is made. Gemini reported confidence 100 in most runs. Its subtle-severity paired coverage (7/73 pairs) is too low for the registered test.


The clean-only quantile scorer (Qwen2.5-0.5B + LoRA, table-held-out) failed its high/low sanity checks and its clean-interval coverage window of [0.80, 0.98] on all six evaluation sets:


\begin{table}[tbp]
\centering
{\small
\setlength{\tabcolsep}{3pt}
\begin{tabular}{lrr}
\toprule
Configuration & Coverage & P+S q / t \\ \midrule
Qwen3-4B / T0 & 0.167 & 15 / 5 \\
Qwen3-4B / T1 & 0.176 & 15 / 5 \\
Qwen3-4B / T1r2 & 0.183 & 14 / 4 \\
Qwen3-4B / T3 & 0.185 & 0 / 0 \\
Qwen3-8B / T0 & 0.176 & 29 / 6 \\
Qwen3-8B / T1 & 0.147 & 29 / 6 \\
\bottomrule
\end{tabular}
}

\caption{Cached learned scorers: all configurations fail eligibility gates. All real-corruption gains are withheld. P+S q / t is the common plausible-plus-subtle query / table count.}
\label{tab:scorer}
\end{table}

\section{Source datasets}

\begin{table}[tbp]
\centering
{\small
\setlength{\tabcolsep}{3pt}
\begin{tabular}{>{\raggedright\arraybackslash}p{1.59in}>{\raggedright\arraybackslash}p{1.53in}}
\toprule
Export / rows / license & Source / SHA-256 prefix \\ \midrule
\path|public_penguins_v3.csv| / 344 / CC0 & Horst et al. palmerpenguins / \path|f204db2c753b| \\
\path|public_auto_mpg_v4.csv| / 398 / CC BY 4.0 & UCI 9, doi:10.24432/C5859H / \path|240f4657d747| \\
\path|public_bike_v4.csv| / 17379 / CC BY 4.0 & UCI 275, doi:10.24432/C5W894 / \path|e03de4ee4ef4| \\
\path|public_wine_quality.csv| / 6497 / CC BY 4.0 & UCI 186, doi:10.24432/C56S3T / \path|428bd8df7313| \\
\path|public_abalone.csv| / 4177 / CC BY 4.0 & UCI 1, doi:10.24432/C55C7W / \path|25f3dc964447| \\
\path|public_student.csv| / 649 / CC BY 4.0 & UCI 320, doi:10.24432/C5TG7T / \path|5be7b02f13f5| \\
\bottomrule
\end{tabular}
}

\caption{Versioned CSV datasets and provenance. SHA-256 prefixes identify the saved exports.}
\label{tab:datasets}
\end{table}

\section{Data, code, and disclosure}

Protocols and dated amendments, raw trajectories, frozen pilot and Experiment A requests with all raw replies, scores, and generation scripts will be released in an anonymized artifact. Palmer Penguins is CC0; the five UCI tables are CC BY 4.0. AI coding assistants helped with code, evidence checks, and drafting; the authors verified all reported numbers against script output.

\end{document}
